\noindent\textbf{(a) Resultado.} \(y=t^2+\cos t-\sin t\).

\noindent\textit{Desarrollo.} Con \(Y=\mathcal L y\), \((s^2+1)Y=s-1+2/s^3+2/s\). Reordenamos
\[Y=2/s^3+s/(s^2+1)-1/(s^2+1).\]
Las inversas son \(t^2\), \(\cos t\) y \(-\sin t\). El término polinómico aporta \((D^2+1)t^2=t^2+2\), los trigonométricos resuelven la homogénea, y los valores en 0 son \(1\) y \(-1\). Los coeficientes continuos garantizan unicidad global.

\medskip

\noindent\textbf{(b) Resultado.} \(y=e^{2t}-3e^{-t}\).

\noindent\textit{Desarrollo.} Los datos deben conservar sus signos al transformar:
\[(s^2-s-2)Y=7-2s,\qquad Y=\frac1{s-2}-\frac3{s+1}.\]
Invertimos para obtener la respuesta. Las raíces \(2\) y \(-1\) anulan el operador, mientras \(y(0)=1-3=-2\), \(y'(0)=2+3=5\). Basta transformar con \(s>2\); la solución en el tiempo existe en todo \(t\ge0\) aunque crezca.

\medskip

\noindent\textbf{(c) Resultado.} \[y=\tfrac15e^{-t}(7\cos t+11\sin t)-\tfrac15(2\cos t-\sin t).\]

\noindent\textit{Desarrollo.} La transformación da \(Y=[s+3+1/(s^2+1)]/(s^2+2s+2)\). Por fracciones parciales,
\[Y=\frac{7(s+1)+11}{5[(s+1)^2+1]}+\frac{-2s+1}{5(s^2+1)}.\]
La inversión produce la respuesta. Como verificación directa, la primera combinación anula el operador \(D^2+2D+2\). Aplicado a \((-2\cos t+\sin t)/5\), da \(\sin t\). Los datos son \((7-2)/5=1\) y \((-7+11+1)/5=1\). La parte transitoria decae y queda una oscilación forzada acotada.
